\section{Conclusion}
\label{sec:conclusion}

A hallucinated tool call can be detected from a single forward pass, and this study is mostly about which signal
to read and when. The first thing to establish is whether an internal judge is needed at all, because that
depends on the model: across the families we tested, either output confidence is uninformative and a trained
probe is the only thing that works, or confidence is already as good as anything we trained, and the two groups
separate cleanly without being explained by size, attention design, release date, or the number of labelled
failures available. Where internals are needed, a probe on the residual stream at the structural positions of
the call is the strongest judge, and attention weights alone come within a few points of it, so an operator
without access to activations still has a usable option.

For attention, one design choice decides everything. Averaging over heads before building the graph cannot
represent the contrast between heads, and for algebraic connectivity it can only overstate how well connected
the graph is, with equality exactly when all heads share a Fiedler vector, a condition every layer we measured
sits far from. Keeping the heads apart moves the attention family from below a length baseline to within
reach of the residual stream, and controls that hold the feature count fixed show that it is the resolution
that does the work. The strongest published attention-only detector already keeps heads apart, and its
features are an in-degree sequence rather than a spectrum, which is why two very different readouts perform
alike once resolution is preserved. The same pattern holds in the residual tier, where the detectors that
read designated positions beat those that pool over the span. Resolution, not the choice of statistic, is what
carries the signal.
